% GENERATED FILE. Edit canonical lessons, then run npm run generate:books.
% canonical-source-hash: fnv1a64:b4ee65786f61919f
% canonical-lessons: TA-C198-pirarttanai-cey, TA-C198-emarru, TA-C198-poy-col, TA-C198-cantai-potu, TA-C198-keli-cey

\chapter{To pray, To cheat, To tell a lie, To quarrel, To tease}
\label{ch:ta-to-pray-to-cheat-to-tell-a-lie-to-quarrel-to-tease}
\hlchaptermodality{198}

\begin{chapteropening}
\textbf{By the end of this chapter:} \emph{I can say to pray, to cheat, to tell a lie, to quarrel and to tease.}

Complete the last lesson of chapter 198: I can say to pray, to cheat, to tell a lie, to quarrel and to tease.
\end{chapteropening}

\section[\texorpdfstring{\ta{பிரார்த்தனை} \ta{செய்} (pirārttaṉai cey)}{பிரார்த்தனை செய் (pirārttaṉai cey)}]{\ta{பிரார்த்தனை} \ta{செய்} (pirārttaṉai cey) — to pray}

\label{lesson:TA-C198-pirarttanai-cey}



\begin{quote}
Before the new one: say the Tamil for to move house, then the Tamil for to contact.
\end{quote}

\subsection*{You'll want to know: \ta{பிரார்த்தனை} \ta{செய்}}
\textbf{\ta{பிரார்த்தனை} \ta{செய்}} — \emph{pirārttaṉai cey} — \textquotedblleft{}to pray\textquotedblright{}.

\begin{grammarlens}[title={What you've built}]
One more word for this chapter.
\end{grammarlens}

\subsection*{Your turn}
\begin{itemize}
\raggedright
  \item \emph{Say it:} \emph{pirārttaṉai cey}
  \item \emph{Say it:} \emph{pirārttaṉai cey}, once more
  \item \emph{Say it:} say \emph{pirārttaṉai cey}
  \item \emph{Recall:} say the Tamil for to move house, then the Tamil for to contact, then say \emph{pirārttaṉai cey} again
\end{itemize}

\begin{tcolorbox}[breakable,colback=teal!4,colframe=teal!35!black,title={Before you move on}]
What does \ta{பிரார்த்தனை} \ta{செய்} mean? (\textquotedblleft{}To pray\textquotedblright{}.) Say it once more.
\end{tcolorbox}
\section[\texorpdfstring{\ta{ஏமாற்று} (ēmāṟṟu)}{ஏமாற்று (ēmāṟṟu)}]{\ta{ஏமாற்று} (ēmāṟṟu) — to cheat, to deceive}

\label{lesson:TA-C198-emarru}



\begin{quote}
Before the new one: say the Tamil for to contact, then the Tamil for to pray.
\end{quote}

\subsection*{You'll want to know: \ta{ஏமாற்று}}
\textbf{\ta{ஏமாற்று}} — \emph{ēmāṟṟu} — \textquotedblleft{}to cheat, to deceive\textquotedblright{}.

\begin{grammarlens}[title={What you've built}]
One more word for this chapter.
\end{grammarlens}

\subsection*{Your turn}
\begin{itemize}
\raggedright
  \item \emph{Say it:} \emph{ēmāṟṟu}
  \item \emph{Say it:} \emph{ēmāṟṟu}, once more
  \item \emph{Say it:} say \emph{ēmāṟṟu}
  \item \emph{Recall:} say the Tamil for to contact, then the Tamil for to pray, then say \emph{ēmāṟṟu} again
\end{itemize}

\begin{tcolorbox}[breakable,colback=teal!4,colframe=teal!35!black,title={Before you move on}]
What does \ta{ஏமாற்று} mean? (\textquotedblleft{}To cheat, to deceive\textquotedblright{}.) Say it once more.
\end{tcolorbox}
\section[\texorpdfstring{\ta{பொய்} \ta{சொல்} (poy col)}{பொய் சொல் (poy col)}]{\ta{பொய்} \ta{சொல்} (poy col) — to tell a lie}

\label{lesson:TA-C198-poy-col}



\begin{quote}
Before the new one: say the Tamil for to pray, then the Tamil for to cheat.
\end{quote}

\subsection*{You'll want to know: \ta{பொய்} \ta{சொல்}}
\textbf{\ta{பொய்} \ta{சொல்}} — \emph{poy col} — \textquotedblleft{}to tell a lie\textquotedblright{}.

\begin{grammarlens}[title={What you've built}]
One more word for this chapter.
\end{grammarlens}

\subsection*{Your turn}
\begin{itemize}
\raggedright
  \item \emph{Say it:} \emph{poy col}
  \item \emph{Say it:} \emph{poy col}, once more
  \item \emph{Say it:} say \emph{poy col}
  \item \emph{Recall:} say the Tamil for to pray, then the Tamil for to cheat, then say \emph{poy col} again
\end{itemize}

\begin{tcolorbox}[breakable,colback=teal!4,colframe=teal!35!black,title={Before you move on}]
What does \ta{பொய்} \ta{சொல்} mean? (\textquotedblleft{}To tell a lie\textquotedblright{}.) Say it once more.
\end{tcolorbox}
\section[\texorpdfstring{\ta{சண்டை} \ta{போடு} (caṇṭai pōṭu)}{சண்டை போடு (caṇṭai pōṭu)}]{\ta{சண்டை} \ta{போடு} (caṇṭai pōṭu) — to quarrel, to fight}

\label{lesson:TA-C198-cantai-potu}



\begin{quote}
Before the new one: say the Tamil for to cheat, then the Tamil for to tell a lie.
\end{quote}

\subsection*{You'll want to know: \ta{சண்டை} \ta{போடு}}
\textbf{\ta{சண்டை} \ta{போடு}} — \emph{caṇṭai pōṭu} — \textquotedblleft{}to quarrel, to fight\textquotedblright{}.

\begin{grammarlens}[title={What you've built}]
One more word for this chapter.
\end{grammarlens}

\subsection*{Your turn}
\begin{itemize}
\raggedright
  \item \emph{Say it:} \emph{caṇṭai pōṭu}
  \item \emph{Say it:} \emph{caṇṭai pōṭu}, once more
  \item \emph{Say it:} say \emph{caṇṭai pōṭu}
  \item \emph{Recall:} say the Tamil for to cheat, then the Tamil for to tell a lie, then say \emph{caṇṭai pōṭu} again
\end{itemize}

\begin{tcolorbox}[breakable,colback=teal!4,colframe=teal!35!black,title={Before you move on}]
What does \ta{சண்டை} \ta{போடு} mean? (\textquotedblleft{}To quarrel, to fight\textquotedblright{}.) Say it once more.
\end{tcolorbox}
\section[\texorpdfstring{\ta{கேலி} \ta{செய்} (kēli cey)}{கேலி செய் (kēli cey)}]{\ta{கேலி} \ta{செய்} (kēli cey) — to tease, to make fun of}

\label{lesson:TA-C198-keli-cey}



\begin{quote}
Before the new one: say the Tamil for to tell a lie, then the Tamil for to quarrel.
\end{quote}

\subsection*{You'll want to know: \ta{கேலி} \ta{செய்}}
\textbf{\ta{கேலி} \ta{செய்}} — \emph{kēli cey} — \textquotedblleft{}to tease, to make fun of\textquotedblright{}.

\begin{grammarlens}[title={What you've built}]
One more word for this chapter.
\end{grammarlens}

\subsection*{Your turn}
\begin{itemize}
\raggedright
  \item \emph{Say it:} \emph{kēli cey}
  \item \emph{Say it:} \emph{kēli cey}, once more
  \item \emph{Say it:} say \emph{kēli cey}
  \item \emph{Recall:} say the Tamil for to tell a lie, then the Tamil for to quarrel, then say \emph{kēli cey} again
\end{itemize}

\begin{tcolorbox}[breakable,colback=teal!4,colframe=teal!35!black,title={Before you move on}]
What does \ta{கேலி} \ta{செய்} mean? (\textquotedblleft{}To tease, to make fun of\textquotedblright{}.) Say it once more.
\end{tcolorbox}
